\documentclass[10pt,a4paper]{report}
\usepackage[utf8]{inputenc}
\usepackage{amsmath}
\usepackage{amsfonts}
\usepackage{amssymb}
\usepackage{amsthm}
\usepackage{hyperref}

\usepackage{multicol}
\usepackage{fancyhdr}
\usepackage[inline]{enumitem}
\usepackage{tikz}
\usepackage{tikz-cd}
\usetikzlibrary{calc}
\usetikzlibrary{shapes.geometric}
\usepackage[margin=0.5in]{geometry}
\usepackage{xcolor}
\usepackage{ esint }

\hypersetup{
    colorlinks=true,
    linkcolor=blue,
    filecolor=magenta,      
    urlcolor=cyan,
    pdftitle={Tensors},
    pdfpagemode=FullScreen,
    }

%\urlstyle{same}

\newcommand{\CLASSNAME}{Math 5230 -- Partial Differential Equations}
\newcommand{\STUDENTNAME}{Paul Carmody}
\newcommand{\ASSIGNMENT}{NOTES }
\newcommand{\DUEDATE}{December 17, 2025}
\newcommand{\SEMESTER}{Fall 2025}
\newcommand{\SCHEDULE}{MW 2:00 -- 3:15}
\newcommand{\ROOM}{Crawford 332}

\newcommand{\MMN}{M_{m\times n}}
\newcommand{\FF}{\mathcal{F}}

\pagestyle{fancy}
\fancyhf{}
\chead{ \fancyplain{}{\CLASSNAME} }
%\chead{ \fancyplain{}{\STUDENTNAME} }
\rhead{\thepage}
\input{../MyMacros}

\newcommand{\RED}[1]{\textcolor{red}{#1}}
\newcommand{\BLUE}[1]{\textcolor{blue}{#1}}

\begin{document}

\begin{center}
	\Large{\CLASSNAME -- \SEMESTER} \\
	\large{ w/Professor Snelson}
\end{center}
\begin{center}
	\STUDENTNAME \\
	\ASSIGNMENT -- \DUEDATE\\
\end{center} 

\begin{description}
	\item \textbf{Transportation Problem and Method Characteristics}
	
	Because the Transportation Equation is
	\begin{align*}
		u_t+b\cdot Du=0
	\end{align*}resembles a linear equation we can use that as inspiration.  Use
	\begin{align*}
		z(s) = u(x+sb, t+s)
	\end{align*}we'll see that $z' = 0$ which is a constant.  Thus, when $s=-t$, intial conditions, we get $z(s) = u(x-tb, 0)$ which is our intial contioins on $u$, namely $g(x)$.  Thus, $u(x,t) = g(x-tb)$
	
	\textbf{Non-Homogenous case}: We find that $z'(s) \ne 0$ and actually is equal to $z'(s) = f(x+sb, t+s$
	\begin{align*}
		u(x,t)-g(x-bt) &= z(0)-z(-t)\\
		&= \int_{-t}^1 z'(s)\, ds \\
		&= \int_0^t f(x+(s-tb,s)\,ds
	\end{align*}thus
	\begin{align*}
		u(x,t) &= g(x-bt) +\int_0^t f(x+(s-tb,s)\,ds
	\end{align*}
	
	\item \textbf{General Laplace solution.}
\begin{remark}[\textbf{Properties of $\Phi$}].

\begin{align*}
	\Phi(x) &= \BINDEF{-\frac{1}{2\pi}\log |x| & (n=2)}{\frac{1}{n(n-2)\alpha(n)}\frac{1}{|x|^{n-2}} & (n\ge 3)} 
\end{align*}
\begin{align}
	D\Phi(y) &= \frac{-1}{n\alpha(n)}\frac{y}{|y|^n},\, (y\ne 0)\\
	\PART{\Phi}{\nu} &= \nu\cdot D\Phi(y)= \frac{-1}{n\alpha(n)\varepsilon^{n-1}},\, y \in \partial B(0,\varepsilon)  \\
	-\Delta\Phi &= \delta_0 \implies \Delta\Phi(x-y)=0,\, y \in U \AND y \ne x
\end{align}implying that there is a point source at 0 or $x$. In (2) $n\alpha(n)\varepsilon^{n-1}$ is the area of the surface of $\partial B(0,\epsilon)$

\end{remark}

	
	\item \textbf{General Poisson solution.}
	
	\begin{align*}
		&\BINDEF{-\Delta u =f, & x \in U}{u=g, & x \in \partial U}\\
		u(x,t) &= \Phi(x) - \phi^x(x) 
	\end{align*}where $\phi^x$ is a Dirichlet Delta function.
	
	\item \textbf{Green's Functions and Dirac Delta functions.}
	
	A Green's Function must satisfy the following:
	\begin{align*}
		G(x,y) &= \BINDEF{ \Delta G(x) = 0, & x \in U}{G(x,y) = \Phi(x-y) + \phi^x(x), & x,y \in \partial U}
	\end{align*}where $\phi^x$ is a Dirichlet Delta function determined by the geometry of $U$.  Then,
	\begin{align*}
		u(x) &= \int_U f(x)G(x,y)\, dy +\int_{\partial U} g(x)\PART{G}{\nu}(x,y) dS
	\end{align*}
	
	\item \textbf{Green's identities.}
\begin{remark}[Green's Identities].
\begin{align*}
\int_U \Delta u\; dx &= \int_{\partial U} \PART{u}{\nu} dS\\
	\int_U Dv\cdot Du\;  dx &= -\int_U u\Delta v \; dx + \int_{\partial U} \PART{v}{\nu}u\;  dS\\
	\int_U u\Delta v - v \Delta u\;  dx &= \int_{\partial U} u\PART{u}{\nu}-v\PART{u}{\nu}\;  dS
\end{align*}
\end{remark}
	
	\item \textbf{Neuman conditions.}
	
	Example from Practice Test 2:
	
	\item \textbf{Energy methods}
	
	Use the functional 
	\begin{align*}
		I[w] &= \int_U \frac{1}{2}|Dw|^2-wf\, dx \\
		w \in \mathcal{A} &= \{ w \in C^2\,|\, w=g \text{ on } \partial U\}
	\end{align*}Dirichlet's Principle then states that
	\begin{align*}
		I[u] &= \min_{w \in \mathcal{A}} I[w] \iff u \text{ is a solution}
	\end{align*}
	
	\item \textbf{Maximum Theorem}
	
	\begin{align*}
		\max_{\bar U} u &= \max_{\partial U} u
	\end{align*}
	
	\item \textbf{Mean Value Theorem}
	
	\begin{align*}
		u(x) &= \fint_U u\, dx = \fint_{\partial U} u \,dS
	\end{align*}
\end{description}

\HLINE
\begin{remark}[Integration by Parts\footnote{see page 9, 9\_17.pdf}].
\begin{align*}
	\int_U f\Delta g\, dx = -\int_U \nabla f \cdot \nabla g \, dx + \int_{\partial U} f \PART{g}{\nu}\, dS
\end{align*}
\begin{align*}
	-\Delta w &= 0 \implies w\Delta w = 0\\
	\int_U w\Delta w\; dx &= -\int_U |Dw|^2\;  dx+\int_{\partial U}wDw\cdot \nu\;  dx = 0\\
	\int_{\partial U} wDw\cdot \nu\;  dx &= \int_U |Dw|^2\;  dx + \int_U w\Delta w\;  \not dx
\end{align*}
\end{remark}

\end{document}